\section{Results on Legal Text}\label{sec:hub:legal}

NQ prompts are short questions.
Two legal datasets repeat the test on legal text, one with prompts of ten words and one with prompts of about \vf{68} words. % src: legal10w piece length; legal128w validation mean 67.6 words (see Datasets paragraph)
A separate hub is trained per dataset, with the architecture and objective of Sections~\ref{sec:hub:design} and~\ref{sec:hub:training}.

\paragraph{Datasets.} Both datasets are cut from the same legal documents, and each prompt starts at a word drawn uniformly over its document. % src: documents/glossary.md (window corpus, piece corpus)
A legal10w prompt holds ten consecutive words.
A legal128w prompt is longer, and its length is limited in tokens.
The prompt grows word by word until it reaches \vf{100} tokens under the tokenizer that yields the most tokens for it. % src: src/v7_8/corpus.py (SEGMENT_MAX_TOKENS 100; tokenizers pythia, olmo, qwen, smollm)
At that limit a prompt averages \vf{68} words on the validation split. % src: artifacts/legal128w/v3/manifests/main-splits.jsonl window_words, validation split (mean 67.6, range 8 to 95)
The \vf{128} in the name is a looser token limit that no prompt reaches, not a word count. % src: src/v7_8/corpus.py (MAX_TOKENS 128); token counts of the first 1,000 validation prompts (maximum 100 under each of the four tokenizers)
Each dataset has \vf{40{,}000} training and \vf{5{,}000} validation prompts. % src: configs/v6_3_legal{128w,10w}.yaml (n_train, n_validation)

\paragraph{Results.} Both legal hubs support the hypothesis (Table~\ref{tab:hub:legal}).
Mean cross-model word-F1 is \vf{0.977} on each dataset, against mismatched-prompt controls of \vf{0.071} on legal10w and \vf{0.152} on legal128w. % src: results/legal{10w,128w}/v6_3/inversion.json cross_mean (hub 0.9767, 0.9769; shuffled 0.0706, 0.1519)
The lowest cross-model direction is Qwen $\to$ OLMo on legal128w, at \vf{0.955} (Figure~\ref{fig:hub:legal-directions}). % src: same files, directions (cross minima 0.9552, 0.9670)
The lower scores into Qwen on NQ do not recur: the cross-model directions into Qwen score \vf{0.975} to \vf{0.986}. % src: same files, olmo->qwen and smollm->qwen (0.9842, 0.9784; 0.9858, 0.9752)
Exact match is higher on the ten-word prompts of legal10w, at \vf{0.83}, than on the long prompts of legal128w, at \vf{0.58}; the comparison does not isolate length from content.

\begin{table}[H]
\caption{Each legal hub reaches a mean cross-model word-F1 of \vf{0.977}, against controls of \vf{0.071} and \vf{0.152}.
Free-running word-F1 and exact match averaged over the six cross-model directions on the \vf{5{,}000} validation prompts of each dataset; each control is evaluated on the same prompts and directions as its score.}
\label{tab:hub:legal}
\centering
\begin{tabular}{lccc}
\toprule
Dataset & Word-F1 & Exact match & Control word-F1 \\
\midrule
legal10w & \vf{0.977} & \vf{0.83} & \vf{0.071} \\ % src: results/legal10w/v6_3/inversion.json cross_mean (0.9767, 0.8291, 0.0706)
legal128w & \vf{0.977} & \vf{0.58} & \vf{0.152} \\ % src: results/legal128w/v6_3/inversion.json cross_mean (0.9769, 0.5814, 0.1519)
\bottomrule
\end{tabular}
\end{table}

\begin{figure}[ht]
\centering
\includegraphics[width=\textwidth]{ch3/direction_recovery_legal}
\caption{Every direction of the two legal hubs reaches a word-F1 of at least 0.95.
Free-running word-F1, with exact match below, by source model (rows) and decoder (columns) on the \vf{5{,}000} validation prompts of each legal dataset; diagonal cells are self-reconstruction.}
\label{fig:hub:legal-directions}
\end{figure}

